\chapter[Unconstrained Optimization \& Gradient Descent]{Unconstrained Optimization, Taylor Expansions and Gradient Descent Convergence}
\label{chap:week08}

% ----------------------------------------------------------------------
% PART 1: EXECUTIVE ROADMAP & THE NARROW RAVINE PARADOX
% ----------------------------------------------------------------------
\section{Executive Roadmap \& Learning Objectives}

The formulation of a machine learning problem as empirical risk minimization culminates in a mathematical optimization problem: finding the parameter vector $\mathbf{w}^* \in \mathbb{R}^d$ that minimizes an unconstrained objective functional $f(\mathbf{w})$. While closed-form solutions exist for Ordinary Least Squares (Chapter~\ref{chap:week03}), modern machine learning models---including logistic regression, support vector machines, and deep neural networks---lack closed-form stationary points. They necessitate iterative numerical optimization. This chapter develops the foundational mathematical theory governing gradient-based and curvature-based algorithms.

\begin{figure}[htbp]
\centering
\begin{tcolorbox}[
    colback=coralbg,
    colframe=coralborder,
    title={\large \textbf{The Narrow Ravine Paradox: Why Steepest Descent Fails to Head for the Minimum}},
    fonttitle=\bfseries\sffamily,
    arc=3mm,
    boxrule=1pt,
    width=\textwidth
]
\textbf{The Intuitive Belief:} Gradient descent follows the direction of \emph{steepest descent} ($-\nabla f(\mathbf{w})$). Intuitively, if you take steps in the direction where the loss drops most steeply, each step must point directly toward the bottom of the bowl (the global minimizer $\mathbf{w}^*$).

\textbf{The Mathematical Reality:} The gradient $\nabla f(\mathbf{w})$ is defined strictly by \emph{infinitesimal local geometry}: it is mathematically orthogonal to the tangent plane of the level set (contour line) passing through $\mathbf{w}$. Unless the loss surface has perfectly spherical contours (equal curvature in all directions), the negative gradient \textbf{almost never points toward the minimum}!

\textbf{The Catastrophe in a Narrow Ravine:} Consider an elongated canyon or ravine where the walls are steep ($L = 100$) but the valley floor slopes gently toward the minimum ($\mu = 1$). If an algorithm is positioned on the canyon wall, the steepest direction does not point along the canyon floor toward the exit; it points directly across the ravine into the opposite wall. As a result, standard gradient descent spends $99.9\%$ of its energy violently ping-ponging back and forth between the cliffs, creeping along the valley floor at a glacial pace of $\mathcal{O}(\kappa)$ steps, where $\kappa = L/\mu$ is the condition number.
\end{tcolorbox}
\caption{The Narrow Ravine Paradox: The fundamental pathology of first-order optimization.}
\label{fig:ravine_paradox}
\end{figure}

\subsection{Prerequisites \& Scaffolding}
Mastery of this chapter builds upon vector calculus, matrix norms, and quadratic form curvature as mapped in Table~\ref{tab:w08_prereqs}.

\begin{table}[htbp]
\centering
\small
\begin{tabularx}{\textwidth}{lXX}
\toprule
\textbf{Foundational Area} & \textbf{Core Concepts Required} & \textbf{Role in Machine Learning Foundations} \\
\midrule
Vector Calculus & Gradient vectors $\nabla f(\mathbf{w})$, Hessian matrices $\nabla^2 f(\mathbf{w})$, Taylor series & Determining local descent directions and second-order curvature \\
Linear Algebra & Spectral radius $\rho(A)$, matrix powers $A^k$, condition number $\kappa$ & Analyzing convergence rates and stability bounds of update operators \\
Mathematical Analysis & Lipschitz continuity, Cauchy sequences, contraction mappings & Establishing formal convergence guarantees for iterative algorithms \\
\bottomrule
\end{tabularx}
\caption{Prerequisite Scaffolding for Chapter 8.}
\label{tab:w08_prereqs}
\end{table}

\subsection{Core Competencies}
By the conclusion of this chapter, students will be able to:
\begin{itemize}
    \item Rigorously state and verify First- and Second-Order Necessary and Sufficient Optimality Conditions.
    \item Formulate $L$-Lipschitz smoothness and prove the fundamental Descent Lemma via integral calculus.
    \item Prove the sublinear $\mathcal{O}(1/k)$ convergence rate of Gradient Descent for convex $L$-smooth functions.
    \item Prove the linear $\mathcal{O}\left(\left(\frac{\kappa - 1}{\kappa + 1}\right)^k\right)$ convergence rate for $\mu$-strongly convex quadratics and determine the optimal step size.
    \item Formulate Newton-Raphson optimization, prove its affine invariance, and analyze the $\mathcal{O}(d^3)$ computational barrier.
    \item Analyze momentum acceleration (Polyak heavy-ball) via discrete dynamical systems.
\end{itemize}

\subsection{Notation \& Dimension Contracts}
Throughout this chapter, we adhere to the standard dimension contracts:
\begin{itemize}
    \item Parameter vector: $\mathbf{w} \in \mathbb{R}^d$ (column vector).
    \item Objective function: $f: \mathbb{R}^d \to \mathbb{R}$.
    \item Gradient vector: $\nabla f(\mathbf{w}) = \left[ \frac{\partial f}{\partial w_1}, \dots, \frac{\partial f}{\partial w_d} \right]^T \in \mathbb{R}^{d \times 1}$.
    \item Hessian matrix: $\nabla^2 f(\mathbf{w}) \in \mathbb{R}^{d \times d}$ where $[\nabla^2 f(\mathbf{w})]_{ij} = \frac{\partial^2 f}{\partial w_i \partial w_j}$.
    \item Learning rate (step size): $\eta > 0$.
    \item Lipschitz smoothness constant: $L = \sup_{\mathbf{w}} \|\nabla^2 f(\mathbf{w})\|_2 = \lambda_{\max}(\nabla^2 f)$.
    \item Strong convexity parameter: $\mu = \inf_{\mathbf{w}} \lambda_{\min}(\nabla^2 f(\mathbf{w})) > 0$.
    \item Condition number: $\kappa \triangleq \frac{L}{\mu} \ge 1$.
\end{itemize}

% ----------------------------------------------------------------------
% PART 2: MATHEMATICAL & COMPUTATIONAL FOUNDATIONS
% ----------------------------------------------------------------------
\section{Mathematical \& Computational Foundations}

\subsection{Multivariable Taylor Approximations}

Let $f: \mathbb{R}^d \to \mathbb{R}$ be twice continuously differentiable ($\mathcal{C}^2$). For any displacement vector $\mathbf{v} \in \mathbb{R}^d$, the scalar function $g(t) = f(\mathbf{w} + t\mathbf{v})$ for $t \in [0, 1]$ allows us to transfer one-dimensional Taylor calculus to $\mathbb{R}^d$:
\begin{align}
g'(t) &= \nabla f(\mathbf{w} + t\mathbf{v})^T \mathbf{v} \\
g''(t) &= \mathbf{v}^T \nabla^2 f(\mathbf{w} + t\mathbf{v}) \mathbf{v}
\end{align}
By Taylor's Theorem with Lagrange remainder on $g(t)$:
\begin{equation}
f(\mathbf{w} + \mathbf{v}) = f(\mathbf{w}) + \nabla f(\mathbf{w})^T \mathbf{v} + \frac{1}{2}\mathbf{v}^T \nabla^2 f(\mathbf{w} + \xi \mathbf{v}) \mathbf{v} \quad \text{for some } \xi \in (0, 1)
\label{eq:taylor_expansion}
\end{equation}
The linear approximation models the function as a hyper-plane, while the quadratic approximation captures local parabolic curvature.

\subsection{Optimality Conditions for Unconstrained Optimization}

\begin{theorem}[First-Order Necessary Condition (FONC)]
Let $f: \mathbb{R}^d \to \mathbb{R}$ be continuously differentiable. If $\mathbf{w}^*$ is a local minimum of $f$, then:
\begin{equation}
\nabla f(\mathbf{w}^*) = \mathbf{0}
\label{eq:fonc}
\end{equation}
A point satisfying $\nabla f(\mathbf{w}) = \mathbf{0}$ is called a \textbf{stationary (critical) point}.
\end{theorem}

\begin{proof}
Suppose $\nabla f(\mathbf{w}^*) \ne \mathbf{0}$. Choose direction $\mathbf{v} = -\nabla f(\mathbf{w}^*)$. By Taylor expansion:
\begin{equation}
f(\mathbf{w}^* + t\mathbf{v}) = f(\mathbf{w}^*) + t \nabla f(\mathbf{w}^*)^T (-\nabla f(\mathbf{w}^*)) + o(t) = f(\mathbf{w}^*) - t \|\nabla f(\mathbf{w}^*)\|_2^2 + o(t)
\end{equation}
For sufficiently small $t > 0$, $-t \|\nabla f(\mathbf{w}^*)\|_2^2 + o(t) < 0$, which implies $f(\mathbf{w}^* + t\mathbf{v}) < f(\mathbf{w}^*)$, directly contradicting that $\mathbf{w}^*$ is a local minimum. Hence $\nabla f(\mathbf{w}^*) = \mathbf{0}$.
\end{proof}

\begin{theorem}[Second-Order Optimality Conditions]
Let $f: \mathbb{R}^d \to \mathbb{R}$ be twice continuously differentiable:
\begin{enumerate}
    \item \textbf{Second-Order Necessary Condition (SONC):} If $\mathbf{w}^*$ is a local minimum, then $\nabla f(\mathbf{w}^*) = \mathbf{0}$ and the Hessian is positive semidefinite:
    \begin{equation}
    \nabla^2 f(\mathbf{w}^*) \succeq 0 \quad (\forall \mathbf{v} \in \mathbb{R}^d, \; \mathbf{v}^T \nabla^2 f(\mathbf{w}^*) \mathbf{v} \ge 0)
    \end{equation}
    \item \textbf{Second-Order Sufficient Condition (SOSC):} If $\nabla f(\mathbf{w}^*) = \mathbf{0}$ and the Hessian is strictly positive definite:
    \begin{equation}
    \nabla^2 f(\mathbf{w}^*) \succ 0 \quad (\forall \mathbf{v} \ne \mathbf{0}, \; \mathbf{v}^T \nabla^2 f(\mathbf{w}^*) \mathbf{v} > 0)
    \end{equation}
    then $\mathbf{w}^*$ is a strict (isolated) local minimum.
\end{enumerate}
\end{theorem}

\begin{ponderbox}[Stationary Points vs.\ Local Minima]
If $\nabla f(\mathbf{w}^*) = \mathbf{0}$ and $\nabla^2 f(\mathbf{w}^*) \succeq 0$, does this guarantee that $\mathbf{w}^*$ is a local minimum?

\textbf{Reflection:} No! Consider the scalar function $f(w) = w^3$. At $w^* = 0$, $f'(0) = 0$ and $f''(0) = 0 \ge 0$. The necessary conditions are fully satisfied. Yet $w^* = 0$ is an inflection point (saddle point), neither a minimum nor a maximum ($f(\epsilon) = \epsilon^3 > 0$ and $f(-\epsilon) = -\epsilon^3 < 0$). Strict positive definiteness ($\nabla^2 f \succ 0$) is required for sufficiency.
\end{ponderbox}

\subsection{Lipschitz Smoothness and the Descent Lemma}

\begin{definitionbox}{$L$-Lipschitz Smoothness}
A continuously differentiable function $f: \mathbb{R}^d \to \mathbb{R}$ is \textbf{$L$-smooth} (its gradient is $L$-Lipschitz continuous) if there exists a constant $L > 0$ such that for all $\mathbf{u}, \mathbf{v} \in \mathbb{R}^d$:
\begin{equation}
\|\nabla f(\mathbf{u}) - \nabla f(\mathbf{v})\|_2 \le L \|\mathbf{u} - \mathbf{v}\|_2
\label{eq:lipschitz_grad}
\end{equation}
For twice-differentiable functions, this condition is mathematically equivalent to bounding the spectral norm of the Hessian everywhere:
\begin{equation}
\|\nabla^2 f(\mathbf{w})\|_2 = \sigma_{\max}(\nabla^2 f(\mathbf{w})) \le L \quad \forall \mathbf{w} \in \mathbb{R}^d
\end{equation}
Geometrically, $L$-smoothness means the curvature of the loss surface cannot exceed $L$ in any direction.
\end{definitionbox}

\begin{theorem}[The Descent Lemma / Quadratic Upper Bound]
If $f$ is $L$-smooth, then for all $\mathbf{u}, \mathbf{v} \in \mathbb{R}^d$:
\begin{equation}
f(\mathbf{v}) \le f(\mathbf{u}) + \nabla f(\mathbf{u})^T (\mathbf{v} - \mathbf{u}) + \frac{L}{2} \|\mathbf{v} - \mathbf{u}\|_2^2
\label{eq:descent_lemma}
\end{equation}
\end{theorem}

\begin{proof}
By the Fundamental Theorem of Calculus along the segment $\mathbf{w}(t) = \mathbf{u} + t(\mathbf{v} - \mathbf{u})$:
\begin{equation}
f(\mathbf{v}) - f(\mathbf{u}) = \int_0^1 \nabla f(\mathbf{u} + t(\mathbf{v} - \mathbf{u}))^T (\mathbf{v} - \mathbf{u}) \, dt
\end{equation}
Add and subtract $\nabla f(\mathbf{u})^T (\mathbf{v} - \mathbf{u})$ inside the integral:
\begin{equation}
f(\mathbf{v}) - f(\mathbf{u}) = \nabla f(\mathbf{u})^T (\mathbf{v} - \mathbf{u}) + \int_0^1 \left( \nabla f(\mathbf{u} + t(\mathbf{v} - \mathbf{u})) - \nabla f(\mathbf{u}) \right)^T (\mathbf{v} - \mathbf{u}) \, dt
\end{equation}
Applying the Cauchy-Schwarz inequality and the $L$-Lipschitz gradient property:
\begin{align}
\left| \int_0^1 \left( \nabla f(\mathbf{u} + t(\mathbf{v} - \mathbf{u})) - \nabla f(\mathbf{u}) \right)^T (\mathbf{v} - \mathbf{u}) \, dt \right|
&\le \int_0^1 \|\nabla f(\mathbf{u} + t(\mathbf{v} - \mathbf{u})) - \nabla f(\mathbf{u})\|_2 \|\mathbf{v} - \mathbf{u}\|_2 \, dt \nonumber \\
&\le \int_0^1 L t \|\mathbf{v} - \mathbf{u}\|_2^2 \, dt = \frac{L}{2} \|\mathbf{v} - \mathbf{u}\|_2^2
\end{align}
Substituting this back establishes the global quadratic upper bound~\eqref{eq:descent_lemma}.
\end{proof}

\begin{theorem}[Sufficient Decrease per Iteration]
Let $f$ be $L$-smooth. Under the Gradient Descent update $\mathbf{w}_{k+1} = \mathbf{w}_k - \eta \nabla f(\mathbf{w}_k)$ with constant step size $\eta \in \left(0, \frac{2}{L}\right)$:
\begin{equation}
f(\mathbf{w}_{k+1}) \le f(\mathbf{w}_k) - \eta \left( 1 - \frac{L\eta}{2} \right) \|\nabla f(\mathbf{w}_k)\|_2^2
\label{eq:sufficient_decrease}
\end{equation}
In particular, for the canonical step size $\eta = \frac{1}{L}$:
\begin{equation}
f(\mathbf{w}_{k+1}) \le f(\mathbf{w}_k) - \frac{1}{2L} \|\nabla f(\mathbf{w}_k)\|_2^2
\label{eq:canonical_decrease}
\end{equation}
\end{theorem}

\begin{proof}
Substitute $\mathbf{u} = \mathbf{w}_k$ and $\mathbf{v} = \mathbf{w}_{k+1} = \mathbf{w}_k - \eta \nabla f(\mathbf{w}_k)$ into the Descent Lemma:
\begin{align}
f(\mathbf{w}_{k+1}) &\le f(\mathbf{w}_k) + \nabla f(\mathbf{w}_k)^T (-\eta \nabla f(\mathbf{w}_k)) + \frac{L}{2} \|-\eta \nabla f(\mathbf{w}_k)\|_2^2 \nonumber \\
&= f(\mathbf{w}_k) - \eta \|\nabla f(\mathbf{w}_k)\|_2^2 + \frac{L\eta^2}{2} \|\nabla f(\mathbf{w}_k)\|_2^2 \nonumber \\
&= f(\mathbf{w}_k) - \eta \left( 1 - \frac{L\eta}{2} \right) \|\nabla f(\mathbf{w}_k)\|_2^2
\end{align}
The coefficient $\eta\left(1 - \frac{L\eta}{2}\right) > 0$ if and only if $0 < \eta < \frac{2}{L}$. Setting the derivative with respect to $\eta$ to zero:
\begin{equation}
\frac{d}{d\eta} \left( \eta - \frac{L\eta^2}{2} \right) = 1 - L\eta = 0 \implies \eta^*_{\text{descent}} = \frac{1}{L}
\end{equation}
Thus, $\eta = \frac{1}{L}$ maximizes the guaranteed one-step decrease: $\Delta f \ge \frac{1}{2L}\|\nabla f(\mathbf{w}_k)\|_2^2$.
\end{proof}

\begin{ponderbox}[The Two Step Sizes: $1/L$ vs.\ $2/(L+\mu)$]
Why do textbooks emphasize two different "optimal" step sizes for Gradient Descent? When should you use $\eta = 1/L$, and when should you use $\eta = \frac{2}{L+\mu}$?

\textbf{Reflection:}
\begin{itemize}
    \item $\eta = 1/L$ is the \textbf{conservative descent step size}. It guarantees that the objective function strictly decreases on \emph{every single step} ($f(\mathbf{w}_{k+1}) \le f(\mathbf{w}_k)$) for \emph{any} $L$-smooth function (even non-convex ones!).
    \item $\eta = \frac{2}{L+\mu}$ is the \textbf{asymptotic contraction step size}. It is designed specifically for \emph{strongly convex quadratics} to minimize the spectral radius of the error transition operator. It yields the fastest possible asymptotic convergence rate $\frac{\kappa-1}{\kappa+1}$, but early steps along the high-curvature axis can overshoot the valley floor (though the overall error in parameter space contractively shrinks).
\end{itemize}
\end{ponderbox}

\subsection{Convergence Rates: Convex vs.\ Strongly Convex}

\begin{theorem}[$\mathcal{O}(1/k)$ Sublinear Convergence for Convex $L$-Smooth Functions]
Let $f: \mathbb{R}^d \to \mathbb{R}$ be convex and $L$-smooth. With constant step size $\eta = \frac{1}{L}$, after $T$ iterations of Gradient Descent:
\begin{equation}
f(\mathbf{w}_T) - f^* \le \frac{2L \|\mathbf{w}_0 - \mathbf{w}^*\|_2^2}{T} = \mathcal{O}\left(\frac{1}{T}\right)
\label{eq:convex_rate}
\end{equation}
To achieve an $\epsilon$-accurate solution ($f(\mathbf{w}_T) - f^* \le \epsilon$), Gradient Descent requires $\mathcal{O}(1/\epsilon)$ iterations.
\end{theorem}

\begin{definitionbox}{Strong Convexity}
A continuously differentiable function $f: \mathbb{R}^d \to \mathbb{R}$ is \textbf{$\mu$-strongly convex} ($\mu > 0$) if for all $\mathbf{u}, \mathbf{v} \in \mathbb{R}^d$:
\begin{equation}
f(\mathbf{v}) \ge f(\mathbf{u}) + \nabla f(\mathbf{u})^T (\mathbf{v} - \mathbf{u}) + \frac{\mu}{2} \|\mathbf{v} - \mathbf{u}\|_2^2
\label{eq:strong_convexity}
\end{equation}
For twice-differentiable functions, this is equivalent to $\nabla^2 f(\mathbf{w}) \succeq \mu I$ for all $\mathbf{w}$. Strong convexity provides a quadratic \emph{lower} bound, guaranteeing that the objective grows quadratically away from the minimum $\mathbf{w}^*$.
\end{definitionbox}

Combining $L$-smoothness (quadratic upper bound) and $\mu$-strong convexity (quadratic lower bound), we obtain the sandwich inequality for all $\mathbf{w}$:
\begin{equation}
\frac{\mu}{2} \|\mathbf{w} - \mathbf{w}^*\|_2^2 \le f(\mathbf{w}) - f^* \le \frac{L}{2} \|\mathbf{w} - \mathbf{w}^*\|_2^2
\label{eq:sandwich_bound}
\end{equation}

% ----------------------------------------------------------------------
% PART 3: VISUAL & GEOMETRIC INTUITION
% ----------------------------------------------------------------------
\section{Visual \& Geometric Intuition: The Geometry of Gradient Descent Trajectories}

Figure~\ref{fig:gd_trajectories} contrasts the convergence trajectory of Gradient Descent on well-conditioned isotropic surfaces against ill-conditioned anisotropic surfaces.

\begin{figure}[htbp]
\centering
\begin{tikzpicture}[
    scale=0.9,
    font=\sffamily\footnotesize
]

% Left: Well-conditioned (circular contours)
\draw[thick, color=primarydark!40] (-4.5, 0) circle (2.2cm);
\draw[thick, color=primarydark!65] (-4.5, 0) circle (1.4cm);
\draw[thick, color=primarydark!90] (-4.5, 0) circle (0.6cm);
\filldraw[color=tealborder] (-4.5, 0) circle (2.5pt) node[below] {$\mathbf{w}^*$};

% Direct descent trajectory
\draw[->, >=stealth, line width=1.4pt, color=tealborder] (-6.3, 1.3) -- (-5.3, 0.6) -- (-4.8, 0.2) -- (-4.5, 0);
\filldraw[color=coralborder] (-6.3, 1.3) circle (2.5pt) node[above left] {$\mathbf{w}_0$};
\node[below, align=center] at (-4.5, -2.5) {\textbf{Isotropic Surface: } $\kappa = 1$ \\ Direct radial descent to $\mathbf{w}^*$};

% Right: Ill-conditioned (stretched elliptical contours)
\begin{scope}[shift={(4.0, 0)}]
    \draw[thick, color=primarydark!40] (0, 0) ellipse (3.6cm and 0.9cm);
    \draw[thick, color=primarydark!65] (0, 0) ellipse (2.4cm and 0.6cm);
    \draw[thick, color=primarydark!90] (0, 0) ellipse (1.2cm and 0.3cm);
    \filldraw[color=tealborder] (0, 0) circle (2.5pt) node[below] {$\mathbf{w}^*$};

    % Zig-zagging trajectory
    \coordinate (W0) at (-3.2, 1.4);
    \coordinate (W1) at (-2.2, -0.7);
    \coordinate (W2) at (-1.3, 0.5);
    \coordinate (W3) at (-0.6, -0.3);
    \coordinate (W4) at (-0.2, 0.1);

    \draw[->, >=stealth, line width=1.4pt, color=traprad] (W0) -- (W1);
    \draw[->, >=stealth, line width=1.4pt, color=traprad] (W1) -- (W2);
    \draw[->, >=stealth, line width=1.4pt, color=traprad] (W2) -- (W3);
    \draw[->, >=stealth, line width=1.4pt, color=traprad] (W3) -- (W4);
    \draw[->, >=stealth, line width=1.4pt, color=traprad] (W4) -- (0, 0);

    % Newton single step arrow
    \draw[->, >=stealth, line width=1.6pt, dashed, color=tealborder] (W0) to[bend left=15] node[above, sloped] {Newton Step (1 iteration)} (0, 0);

    \filldraw[color=coralborder] (W0) circle (2.5pt) node[above] {$\mathbf{w}_0$};
    \node[below, align=center] at (0, -2.5) {\textbf{Ill-Conditioned Surface: } $\kappa \gg 1$ \\ Severe cross-ravine zig-zagging};
\end{scope}

\end{tikzpicture}
\caption{The Geometry of Optimization Trajectories: On isotropic surfaces ($\kappa = 1$), gradient arrows point straight to the minimizer $\mathbf{w}^*$. On ill-conditioned surfaces ($\kappa \gg 1$), gradients point nearly perpendicular to the valley axis, causing severe oscillation. In contrast, Newton's method (dashed line) jumps directly to $\mathbf{w}^*$ in a single iteration.}
\label{fig:gd_trajectories}
\end{figure}

% ----------------------------------------------------------------------
% PART 4: DEEP INTUITION & CONCEPTUAL DERIVATIONS
% ----------------------------------------------------------------------
\section{Deep Intuition \& Conceptual Derivations}

\subsection{Exact Dynamics on Strictly Convex Quadratics}
Consider minimizing the quadratic objective:
\begin{equation}
f(\mathbf{w}) = \frac{1}{2}\mathbf{w}^T A \mathbf{w} - \mathbf{b}^T \mathbf{w}
\label{eq:quadratic_obj}
\end{equation}
where $A \in \mathbb{R}^{d \times d}$ is symmetric positive definite ($A \succ 0$) with eigenvalues:
\begin{equation}
0 < \mu = \lambda_{\min}(A) = \lambda_1 \le \lambda_2 \le \dots \le \lambda_d = \lambda_{\max}(A) = L
\end{equation}
The gradient is $\nabla f(\mathbf{w}) = A\mathbf{w} - \mathbf{b}$, and the unique global minimizer is $\mathbf{w}^* = A^{-1}\mathbf{b}$.

The Gradient Descent recurrence with step size $\eta$ is:
\begin{equation}
\mathbf{w}_{k+1} = \mathbf{w}_k - \eta (A\mathbf{w}_k - \mathbf{b})
\end{equation}
Subtracting $\mathbf{w}^* = A^{-1}\mathbf{b}$ from both sides:
\begin{align}
\mathbf{w}_{k+1} - \mathbf{w}^* &= \mathbf{w}_k - \mathbf{w}^* - \eta (A\mathbf{w}_k - A\mathbf{w}^*) \nonumber \\
&= (I - \eta A)(\mathbf{w}_k - \mathbf{w}^*)
\end{align}
Iterating this recurrence $k$ times:
\begin{equation}
\mathbf{w}_k - \mathbf{w}^* = (I - \eta A)^k (\mathbf{w}_0 - \mathbf{w}^*)
\label{eq:gd_matrix_power}
\end{equation}

\subsubsection{Eigen-Decomposition of the Error Vector}
Because $A$ is symmetric, it admits an orthonormal spectral decomposition $A = Q \Lambda Q^T = \sum_{i=1}^d \lambda_i \mathbf{q}_i \mathbf{q}_i^T$.
Expressing the initial error in the eigenbasis of $A$:
\begin{equation}
\mathbf{w}_0 - \mathbf{w}^* = \sum_{i=1}^d \alpha_i \mathbf{q}_i, \quad \text{where } \alpha_i = \mathbf{q}_i^T (\mathbf{w}_0 - \mathbf{w}^*)
\end{equation}
Since $(I - \eta A) \mathbf{q}_i = (1 - \eta \lambda_i) \mathbf{q}_i$, the error after $k$ steps uncouples into $d$ independent scalar progressions:
\begin{equation}
\mathbf{w}_k - \mathbf{w}^* = \sum_{i=1}^d (1 - \eta \lambda_i)^k \alpha_i \mathbf{q}_i
\label{eq:decoupled_error}
\end{equation}

\subsubsection{Stability Condition}
Convergence requires that every modal error shrinks to zero:
\begin{equation}
\lim_{k \to \infty} (1 - \eta \lambda_i)^k = 0 \iff |1 - \eta \lambda_i| < 1 \quad \forall i \in \{1, \dots, d\}
\end{equation}
This holds if and only if $-1 < 1 - \eta \lambda_i < 1$, which yields:
\begin{equation}
0 < \eta < \frac{2}{\lambda_{\max}} = \frac{2}{L}
\label{eq:stability_threshold}
\end{equation}
If $\eta \ge \frac{2}{L}$, the error along the eigenvector $\mathbf{q}_d$ of maximum eigenvalue $\lambda_{\max} = L$ oscillates with factor $|1 - \eta L| \ge 1$ and fails to converge (exploding if $\eta > 2/L$).

\subsubsection{Optimal Learning Rate and Convergence Factor}
The asymptotic convergence rate is governed by the spectral radius $\rho(I - \eta A) = \max_i |1 - \eta \lambda_i| = \max(|1 - \eta \mu|, |1 - \eta L|)$.
To minimize this worst-case contraction factor:
\begin{equation}
1 - \eta^* \mu = -(1 - \eta^* L) \implies \eta^* = \frac{2}{L + \mu}
\label{eq:optimal_eta}
\end{equation}
Substituting $\eta^*$ back into the contraction factor:
\begin{equation}
\rho^* = 1 - \eta^* \mu = 1 - \frac{2\mu}{L + \mu} = \frac{L - \mu}{L + \mu} = \frac{\frac{L}{\mu} - 1}{\frac{L}{\mu} + 1} = \frac{\kappa - 1}{\kappa + 1}
\label{eq:condition_contraction}
\end{equation}
Taking the Euclidean norm of~\eqref{eq:decoupled_error}:
\begin{equation}
\|\mathbf{w}_k - \mathbf{w}^*\|_2 \le \left( \frac{\kappa - 1}{\kappa + 1} \right)^k \|\mathbf{w}_0 - \mathbf{w}^*\|_2
\label{eq:linear_conv_rate}
\end{equation}
This is \textbf{linear convergence} (geometric decay in error). When the condition number is large ($\kappa \gg 1$), a first-order Taylor expansion reveals:
\begin{equation}
\frac{\kappa - 1}{\kappa + 1} = \frac{1 - 1/\kappa}{1 + 1/\kappa} \approx \left(1 - \frac{1}{\kappa}\right)\left(1 - \frac{1}{\kappa}\right) \approx 1 - \frac{2}{\kappa}
\end{equation}
To reduce the initial error by a factor of $\epsilon$, the number of iterations required is:
\begin{equation}
\left( 1 - \frac{2}{\kappa} \right)^k \le \epsilon \implies k \approx \frac{\kappa}{2} \ln\left(\frac{1}{\epsilon}\right) = \mathcal{O}(\kappa \ln(1/\epsilon))
\end{equation}
If $\kappa = 10{,}000$, gradient descent requires over $50{,}000$ iterations to achieve moderate precision!

\subsection{Newton-Raphson Method: Curvature Inversion}
Newton's method addresses the condition number curse by modeling the loss surface with its full second-order Taylor expansion around $\mathbf{w}_k$:
\begin{equation}
q(\mathbf{w}) = f(\mathbf{w}_k) + \nabla f(\mathbf{w}_k)^T (\mathbf{w} - \mathbf{w}_k) + \frac{1}{2}(\mathbf{w} - \mathbf{w}_k)^T \nabla^2 f(\mathbf{w}_k) (\mathbf{w} - \mathbf{w}_k)
\end{equation}
Minimizing this quadratic model with respect to $\mathbf{w}$ yields:
\begin{equation}
\nabla q(\mathbf{w}) = \nabla f(\mathbf{w}_k) + \nabla^2 f(\mathbf{w}_k)(\mathbf{w} - \mathbf{w}_k) = \mathbf{0}
\end{equation}
Assuming $\nabla^2 f(\mathbf{w}_k) \succ 0$, the exact minimizer of the local quadratic model defines the Newton update:
\begin{equation}
\mathbf{w}_{k+1} = \mathbf{w}_k - [\nabla^2 f(\mathbf{w}_k)]^{-1} \nabla f(\mathbf{w}_k)
\label{eq:newton_update}
\end{equation}

\subsubsection{One-Step Convergence on Quadratics}
For any quadratic $f(\mathbf{w}) = \frac{1}{2}\mathbf{w}^T A \mathbf{w} - \mathbf{b}^T \mathbf{w}$ with $A \succ 0$:
The gradient is $\nabla f(\mathbf{w}) = A\mathbf{w} - \mathbf{b}$ and the Hessian is $\nabla^2 f(\mathbf{w}) = A$.
Applying Newton's update from \emph{any} initial point $\mathbf{w}_0$:
\begin{equation}
\mathbf{w}_1 = \mathbf{w}_0 - A^{-1}(A\mathbf{w}_0 - \mathbf{b}) = \mathbf{w}_0 - \mathbf{w}_0 + A^{-1}\mathbf{b} = A^{-1}\mathbf{b} = \mathbf{w}^*
\end{equation}
Newton's method finds the exact global minimum in \textbf{precisely one iteration}, completely independent of the condition number $\kappa$!

\subsubsection{Affine Invariance of Newton's Method}
Unlike Gradient Descent, Newton's method is completely invariant to non-singular linear changes of coordinates.
Let $\mathbf{y} = B \mathbf{w}$ where $B \in \mathbb{R}^{d \times d}$ is invertible. Define $\tilde{f}(\mathbf{y}) = f(B^{-1}\mathbf{y})$.
By the chain rule:
\begin{align}
\nabla \tilde{f}(\mathbf{y}) &= B^{-T} \nabla f(\mathbf{w}) \\
\nabla^2 \tilde{f}(\mathbf{y}) &= B^{-T} \nabla^2 f(\mathbf{w}) B^{-1}
\end{align}
The Newton step in $\mathbf{y}$-coordinates is:
\begin{align}
\Delta \mathbf{y} &= -[\nabla^2 \tilde{f}(\mathbf{y})]^{-1} \nabla \tilde{f}(\mathbf{y}) \nonumber \\
&= -\left[ B^{-T} \nabla^2 f(\mathbf{w}) B^{-1} \right]^{-1} \left[ B^{-T} \nabla f(\mathbf{w}) \right] \nonumber \\
&= -B [\nabla^2 f(\mathbf{w})]^{-1} B^T B^{-T} \nabla f(\mathbf{w}) \nonumber \\
&= B \left( -[\nabla^2 f(\mathbf{w})]^{-1} \nabla f(\mathbf{w}) \right) = B \Delta \mathbf{w}
\end{align}
The update in the transformed space is identical to the transformation of the update in the original space: $\mathbf{y}_{k+1} = B \mathbf{w}_{k+1}$. Newton's method effectively rescales the coordinate axes so that the contours appear perfectly spherical.

\begin{ponderbox}[Why isn't Newton's Method used for Deep Learning?]
If Newton's method converges quadratically ($\| \mathbf{w}_{k+1} - \mathbf{w}^* \| \le C \| \mathbf{w}_k - \mathbf{w}^* \|^2$) and is invariant to ill-conditioning, why do modern deep learning frameworks (PyTorch, TensorFlow) almost universally rely on first-order variants (SGD, Adam)?

\textbf{Reflection:} Three severe obstacles prevent direct application of Newton's method in deep learning:
\begin{enumerate}
    \item \textbf{Memory Storage $\mathcal{O}(d^2)$:} A modest modern neural network has $d = 10^8$ parameters. Storing the Hessian requires $d^2 = 10^{16}$ floats ($40$ Petabytes of RAM).
    \item \textbf{Computational Inversion $\mathcal{O}(d^3)$:} Inverting a $10^8 \times 10^8$ matrix takes $\approx 10^{24}$ FLOPs, requiring centuries of computation per gradient step.
    \item \textbf{Non-Convexity and Saddle Points:} In high dimensions, random stationary points are overwhelmingly saddle points rather than local minima. Standard Newton's method is attracted to local \emph{maxima} and saddle points because it only seeks $\nabla f = \mathbf{0}$, jumping violently in directions of negative curvature ($\lambda_i < 0$).
\end{enumerate}
\end{ponderbox}

\subsection{Polyak Momentum: Heavy-Ball Acceleration}
The Heavy-Ball momentum method introduces a physical inertia term:
\begin{equation}
\mathbf{w}_{k+1} = \mathbf{w}_k - \alpha \nabla f(\mathbf{w}_k) + \beta (\mathbf{w}_k - \mathbf{w}_{k-1})
\label{eq:momentum_update}
\end{equation}
where $\beta \in [0, 1)$ is the momentum parameter.

\subsubsection{Discrete Dynamical System Formulation}
On a one-dimensional quadratic $f(w) = \frac{1}{2} a w^2$ with $a > 0$:
\begin{equation}
w_{k+1} = (1 - \alpha a + \beta) w_k - \beta w_{k-1}
\end{equation}
Stacking the state vector $\mathbf{u}_k = [w_k, w_{k-1}]^T \in \mathbb{R}^2$:
\begin{equation}
\begin{bmatrix} w_{k+1} \\ w_k \end{bmatrix} = \begin{bmatrix} 1 - \alpha a + \beta & -\beta \\ 1 & 0 \end{bmatrix} \begin{bmatrix} w_k \\ w_{k-1} \end{bmatrix} \implies \mathbf{u}_{k+1} = T \mathbf{u}_k
\label{eq:momentum_system}
\end{equation}
The characteristic polynomial of the transition matrix $T$ is:
\begin{equation}
\det(T - \lambda I) = \lambda^2 - (1 - \alpha a + \beta)\lambda + \beta = 0
\end{equation}
By Vieta's formulas, the product of the eigenvalues is $\lambda_1 \lambda_2 = \det(T) = \beta$.
When the roots are complex conjugates, their magnitudes are identically $|\lambda_1| = |\lambda_2| = \sqrt{\beta}$!
By optimally setting:
\begin{equation}
\alpha^* = \frac{4}{(\sqrt{L} + \sqrt{\mu})^2}, \quad \beta^* = \left( \frac{\sqrt{L} - \sqrt{\mu}}{\sqrt{L} + \sqrt{\mu}} \right)^2 = \left( \frac{\sqrt{\kappa} - 1}{\sqrt{\kappa} + 1} \right)^2
\label{eq:optimal_momentum}
\end{equation}
the spectral radius of the transition matrix becomes:
\begin{equation}
\rho(T) = \sqrt{\beta^*} = \frac{\sqrt{\kappa} - 1}{\sqrt{\kappa} + 1}
\end{equation}
This improves the condition number dependence from $\mathcal{O}(\kappa)$ to $\mathcal{O}(\sqrt{\kappa})$. For $\kappa = 10{,}000$, $\kappa$ drops from $10{,}000$ to $\sqrt{\kappa} = 100$, accelerating convergence by a factor of $100\times$!

% ----------------------------------------------------------------------
% PART 5: MULTI-PAGE WORKED EXAMPLES
% ----------------------------------------------------------------------
\section{Multi-Page Breakable Numerical Worked Examples}

\begin{examplebox}[breakable]{Step-by-Step Gradient Descent Dynamics on an Ill-Conditioned Quadratic}
\textbf{Problem Statement:}
Consider the quadratic objective function:
\begin{equation}
f(x_1, x_2) = \frac{1}{2} x_1^2 + 5 x_2^2 = \frac{1}{2} \mathbf{x}^T \begin{bmatrix} 1 & 0 \\ 0 & 10 \end{bmatrix} \mathbf{x}
\end{equation}
The unique global minimum is at $\mathbf{x}^* = [0, 0]^T$ with $f^* = 0$.
The starting point is $\mathbf{x}_0 = [10.0, 1.0]^T$.

\textbf{Task Requirements:}
\begin{enumerate}
    \item Identify the eigenvalues $\lambda_1, \lambda_2$, Lipschitz constant $L$, strong convexity parameter $\mu$, and condition number $\kappa$.
    \item Determine the maximum non-divergent step size $\eta_{\max}$ and the theoretical optimal step size $\eta^*$.
    \item Manually calculate the exact values of iterates $\mathbf{x}_1, \mathbf{x}_2$ and objective losses $f(\mathbf{x}_0), f(\mathbf{x}_1), f(\mathbf{x}_2)$ under the optimal learning rate $\eta^* = \frac{2}{L+\mu}$.
    \item Contrast this with the divergent behavior under step size $\eta = 0.22 > \eta_{\max}$.
\end{enumerate}

\tcblower

\textbf{Step 1: Curvature Parameters and Condition Number}
The Hessian is constant:
\begin{equation}
\nabla^2 f(\mathbf{x}) = A = \begin{bmatrix} 1 & 0 \\ 0 & 10 \end{bmatrix}
\end{equation}
The eigenvalues are:
\begin{equation}
\lambda_1 = \mu = 1.0, \quad \lambda_2 = L = 10.0
\end{equation}
The condition number is:
\begin{equation}
\kappa = \frac{L}{\mu} = \frac{10.0}{1.0} = 10
\end{equation}

\textbf{Step 2: Step Size Boundaries}
\begin{itemize}
    \item Maximum stable step size:
    \begin{equation}
    \eta_{\max} = \frac{2}{L} = \frac{2}{10.0} = 0.200
    \end{equation}
    \item Optimal asymptotic step size:
    \begin{equation}
    \eta^* = \frac{2}{L + \mu} = \frac{2}{10 + 1} = \frac{2}{11} \approx 0.181818
    \end{equation}
\end{itemize}

\textbf{Step 3: Exact Iterate Computations Under $\eta^* = \frac{2}{11}$}
The gradient is:
\begin{equation}
\nabla f(\mathbf{x}) = \begin{bmatrix} x_1 \\ 10 x_2 \end{bmatrix}
\end{equation}
The initial objective value at $\mathbf{x}_0 = [10, 1]^T$ is:
\begin{equation}
f(\mathbf{x}_0) = \frac{1}{2}(10)^2 + 5(1)^2 = 50 + 5 = 55.000
\end{equation}
The gradient at $\mathbf{x}_0$ is:
\begin{equation}
\nabla f(\mathbf{x}_0) = \begin{bmatrix} 10.0 \\ 10(1.0) \end{bmatrix} = \begin{bmatrix} 10.0 \\ 10.0 \end{bmatrix}
\end{equation}

\textbf{Iteration 1 ($\mathbf{x}_1$):}
\begin{align}
\mathbf{x}_1 &= \mathbf{x}_0 - \eta^* \nabla f(\mathbf{x}_0) \nonumber \\
&= \begin{bmatrix} 10 \\ 1 \end{bmatrix} - \frac{2}{11} \begin{bmatrix} 10 \\ 10 \end{bmatrix}
= \begin{bmatrix} 10 - \frac{20}{11} \\ 1 - \frac{20}{11} \end{bmatrix}
= \begin{bmatrix} \frac{90}{11} \\ -\frac{9}{11} \end{bmatrix} \approx \begin{bmatrix} 8.1818 \\ -0.8182 \end{bmatrix}
\end{align}
Notice the critical diagnostic event: $x_2$ has flipped sign from $+1.0$ to $-0.8182$! It overshot the bottom of the ravine and landed on the opposite wall.
The objective value at $\mathbf{x}_1$ is:
\begin{align}
f(\mathbf{x}_1) &= \frac{1}{2}\left(\frac{90}{11}\right)^2 + 5\left(-\frac{9}{11}\right)^2 \nonumber \\
&= \frac{1}{2} \cdot \frac{8100}{121} + 5 \cdot \frac{81}{121} = \frac{4050 + 405}{121} = \frac{4455}{121} \approx 36.8182
\end{align}

\textbf{Iteration 2 ($\mathbf{x}_2$):}
The gradient at $\mathbf{x}_1$ is:
\begin{equation}
\nabla f(\mathbf{x}_1) = \begin{bmatrix} \frac{90}{11} \\ 10\left(-\frac{9}{11}\right) \end{bmatrix} = \begin{bmatrix} \frac{90}{11} \\ -\frac{90}{11} \end{bmatrix}
\end{equation}
Now update:
\begin{align}
\mathbf{x}_2 &= \mathbf{x}_1 - \eta^* \nabla f(\mathbf{x}_1) \nonumber \\
&= \begin{bmatrix} \frac{90}{11} \\ -\frac{9}{11} \end{bmatrix} - \frac{2}{11} \begin{bmatrix} \frac{90}{11} \\ -\frac{90}{11} \end{bmatrix}
= \begin{bmatrix} \frac{90}{11}\left(1 - \frac{2}{11}\right) \\ -\frac{9}{11}\left(1 - \frac{20}{11}\right) \end{bmatrix}
= \begin{bmatrix} \frac{90}{11} \cdot \frac{9}{11} \\ -\frac{9}{11} \left(-\frac{9}{11}\right) \end{bmatrix}
= \begin{bmatrix} \frac{810}{121} \\ \frac{81}{121} \end{bmatrix} \approx \begin{bmatrix} 6.6942 \\ +0.6694 \end{bmatrix}
\end{align}
The sign of $x_2$ flips back to positive ($+0.6694$).
The coordinate progression for any step $k$ is given analytically by:
\begin{align}
x_1^{(k)} &= \left(1 - \eta^* \lambda_1\right)^k x_1^{(0)} = \left(1 - \frac{2}{11}\right)^k (10) = 10 \left(\frac{9}{11}\right)^k \\
x_2^{(k)} &= \left(1 - \eta^* \lambda_2\right)^k x_2^{(0)} = \left(1 - \frac{20}{11}\right)^k (1) = 1 \left(-\frac{9}{11}\right)^k
\end{align}
Observe that both coordinate errors contract at the exact same rate:
\begin{equation}
|x_1^{(k)}| = 10 \left(\frac{9}{11}\right)^k, \quad |x_2^{(k)}| = 1 \left(\frac{9}{11}\right)^k
\end{equation}
The contraction factor is precisely $\rho^* = \frac{\kappa - 1}{\kappa + 1} = \frac{10 - 1}{10 + 1} = \frac{9}{11} \approx 0.8182$.
The objective value at $\mathbf{x}_2$ is:
\begin{equation}
f(\mathbf{x}_2) = \left(\frac{9}{11}\right)^2 f(\mathbf{x}_1) = \frac{81}{121} \times \frac{4455}{121} \approx 24.646
\end{equation}

\textbf{Step 4: Divergence with $\eta = 0.22 > \eta_{\max}$}
Suppose a user chooses step size $\eta = 0.22 > 0.20$.
The contraction factors along the two axes are:
\begin{align}
1 - \eta \lambda_1 &= 1 - 0.22(1) = +0.78 \quad (\text{stable, contracts}) \\
1 - \eta \lambda_2 &= 1 - 0.22(10) = 1 - 2.20 = -1.20 \quad (\text{UNSTABLE!})
\end{align}
Along the $x_2$ axis, the error expands exponentially:
\begin{equation}
x_2^{(k)} = (-1.20)^k \cdot (1.0)
\end{equation}
\begin{itemize}
    \item $x_2^{(0)} = +1.0$
    \item $x_2^{(1)} = -1.20$
    \item $x_2^{(2)} = +1.44$
    \item $x_2^{(5)} = -2.488$
    \item $x_2^{(20)} \approx +38.33$
    \item $x_2^{(100)} \approx 8.28 \times 10^7$
\end{itemize}
Even though $x_1 \to 0$, $x_2$ explodes violently to $\pm \infty$, triggering numerical overflow (\texttt{NaN} in Python).
\end{examplebox}

\begin{examplebox}[breakable]{Newton's Method vs.\ Gradient Descent on a Multi-Variable Objective}
\textbf{Problem Statement:}
Consider the non-quadratic function:
\begin{equation}
f(x_1, x_2) = x_1^4 + 2x_1 x_2 + 2x_2^2
\end{equation}
Starting at the initial test point $\mathbf{x}_0 = [1.0, 1.0]^T$:
\begin{enumerate}
    \item Compute the gradient $\nabla f(\mathbf{x}_0)$.
    \item Compute the Hessian matrix $\nabla^2 f(\mathbf{x}_0)$ and verify that it is strictly positive definite.
    \item Compute the exact inverse Hessian $[\nabla^2 f(\mathbf{x}_0)]^{-1}$.
    \item Compute the first Newton iterate $\mathbf{x}_1^{\text{Newton}}$ and compare it with the Gradient Descent iterate $\mathbf{x}_1^{\text{GD}}$ using $\eta = 0.1$.
\end{enumerate}

\tcblower

\textbf{Step 1: Gradient Vector}
The partial derivatives are:
\begin{align}
\frac{\partial f}{\partial x_1} &= 4x_1^3 + 2x_2 \\
\frac{\partial f}{\partial x_2} &= 2x_1 + 4x_2
\end{align}
Evaluating at $\mathbf{x}_0 = [1, 1]^T$:
\begin{equation}
\nabla f(\mathbf{x}_0) = \begin{bmatrix} 4(1)^3 + 2(1) \\ 2(1) + 4(1) \end{bmatrix} = \begin{bmatrix} 6 \\ 6 \end{bmatrix}
\end{equation}

\textbf{Step 2: Hessian Matrix Evaluation}
The second partial derivatives are:
\begin{align}
\frac{\partial^2 f}{\partial x_1^2} &= 12x_1^2, \quad \frac{\partial^2 f}{\partial x_1 \partial x_2} = 2 \\
\frac{\partial^2 f}{\partial x_2 \partial x_1} &= 2, \quad \frac{\partial^2 f}{\partial x_2^2} = 4
\end{align}
Evaluating at $\mathbf{x}_0 = [1, 1]^T$:
\begin{equation}
H_0 = \nabla^2 f(\mathbf{x}_0) = \begin{bmatrix} 12 & 2 \\ 2 & 4 \end{bmatrix}
\end{equation}
Check positive definiteness:
\begin{itemize}
    \item Leading principal minor 1: $\Delta_1 = 12 > 0$.
    \item Leading principal minor 2: $\Delta_2 = \det(H_0) = (12)(4) - (2)(2) = 48 - 4 = 44 > 0$.
\end{itemize}
By Sylvester's Criterion, $H_0 \succ 0$ strictly.

\textbf{Step 3: Inverse Hessian Matrix}
Using the $2 \times 2$ matrix inverse formula:
\begin{equation}
H_0^{-1} = \frac{1}{\det(H_0)} \begin{bmatrix} 4 & -2 \\ -2 & 12 \end{bmatrix} = \frac{1}{44} \begin{bmatrix} 4 & -2 \\ -2 & 12 \end{bmatrix} = \begin{bmatrix} \frac{1}{11} & -\frac{1}{22} \\ -\frac{1}{22} & \frac{3}{11} \end{bmatrix}
\end{equation}

\textbf{Step 4: Comparison of Iterates}
\begin{itemize}
    \item \textbf{Newton Iteration:}
    \begin{align}
    \Delta \mathbf{x}^{\text{Newton}} &= -H_0^{-1} \nabla f(\mathbf{x}_0) = -\frac{1}{44} \begin{bmatrix} 4 & -2 \\ -2 & 12 \end{bmatrix} \begin{bmatrix} 6 \\ 6 \end{bmatrix} \nonumber \\
    &= -\frac{1}{44} \begin{bmatrix} 24 - 12 \\ -12 + 72 \end{bmatrix} = -\frac{1}{44} \begin{bmatrix} 12 \\ 60 \end{bmatrix} = \begin{bmatrix} -\frac{3}{11} \\ -\frac{15}{11} \end{bmatrix} \approx \begin{bmatrix} -0.2727 \\ -1.3636 \end{bmatrix}
    \end{align}
    The updated parameter vector is:
    \begin{equation}
    \mathbf{x}_1^{\text{Newton}} = \mathbf{x}_0 + \Delta \mathbf{x}^{\text{Newton}} = \begin{bmatrix} 1 \\ 1 \end{bmatrix} + \begin{bmatrix} -\frac{3}{11} \\ -\frac{15}{11} \end{bmatrix} = \begin{bmatrix} \frac{8}{11} \\ -\frac{4}{11} \end{bmatrix} \approx \begin{bmatrix} 0.7273 \\ -0.3636 \end{bmatrix}
    \end{equation}

    \item \textbf{Gradient Descent Iteration ($\eta = 0.1$):}
    \begin{equation}
    \mathbf{x}_1^{\text{GD}} = \mathbf{x}_0 - \eta \nabla f(\mathbf{x}_0) = \begin{bmatrix} 1 \\ 1 \end{bmatrix} - 0.1 \begin{bmatrix} 6 \\ 6 \end{bmatrix} = \begin{bmatrix} 0.4 \\ 0.4 \end{bmatrix}
    \end{equation}
\end{itemize}
Notice that Gradient Descent moves equally along both coordinates ($0.4, 0.4$) because the gradient vector had identical components ($6, 6$). However, Newton's method recognized that the curvature along $x_1$ ($12$) is three times greater than along $x_2$ ($4$), so it took a modest step along $x_1$ and a much larger, corrective step along $x_2$, adapting automatically to the local geometry.
\end{examplebox}

% ----------------------------------------------------------------------
% PART 6: EXPLICIT GOTCHAS & DIAGNOSTIC MISCONCEPTIONS
% ----------------------------------------------------------------------
\section{Explicit Gotchas \& Diagnostic Misconceptions}

\begin{trapbox}{Gotcha 1: Selecting Learning Rate Above the Theoretical Stability Bound ($2/L$)}
A student sets learning rate $\eta = 0.25$ on a problem where $L = 10$. They observe the loss oscillating wildly and overflowing to \texttt{NaN} within 15 iterations.
\textbf{The Trap:} The divergence boundary is strictly $\eta_{\max} = \frac{2}{L}$. For $L = 10$, $\eta_{\max} = \frac{2}{10} = 0.20$. Any step size $\eta \ge 0.20$ creates an error amplification factor $|1 - \eta L| \ge 1$. The error grows exponentially as $|1 - \eta L|^k \to \infty$. Step sizes must satisfy $\eta < \frac{2}{L}$. In machine learning practice, setting $\eta \le \frac{1}{L}$ guarantees monotonic objective decrease.
\end{trapbox}

\begin{trapbox}{Gotcha 2: Assuming Gradient Descent Always Finds the Global Minimum}
A practitioner executes gradient descent on a complex deep neural network and asserts the final parameter vector is the global minimizer.
\textbf{The Trap:} Gradient descent is guaranteed to find the global minimum only if the objective function is \textbf{convex}. On non-convex loss surfaces, gradient descent can be trapped in suboptimal local minima or stall indefinitely near saddle points and plateaus where $\nabla f(\mathbf{w}) \approx \mathbf{0}$.
\end{trapbox}

\begin{trapbox}{Gotcha 3: Conflating Stationary Points with Local Minima}
An optimization script terminates when $\|\nabla f(\mathbf{w})\|_2 < 10^{-5}$. The engineer logs that a local minimum has been discovered.
\textbf{The Trap:} A vanishing gradient ($\nabla f = \mathbf{0}$) is a \textbf{necessary} condition, not a sufficient condition. On non-convex surfaces, stationary points can be local maxima or saddle points. Verifying a local minimum requires evaluating the Hessian eigenvalues: $\lambda_i(\nabla^2 f) > 0$ for all $i$.
\end{trapbox}

\begin{trapbox}{Gotcha 4: Inverting Hessians in High Dimensions (The $\mathcal{O}(d^3)$ Barrier)}
An engineer decides to replace gradient descent with Newton's method on a modern computer vision model with $d = 10^7$ parameters.
\textbf{The Trap:} Storing the Hessian requires storing $d^2 = 10^{14}$ floating-point values ($800$ Terabytes of RAM). Inverting it via Gaussian elimination requires $\mathcal{O}(d^3) = 10^{21}$ FLOPs, which would take centuries of compute time. Second-order methods in modern ML require quasi-Newton approximations (L-BFGS) or Hessian-free conjugate gradient methods.
\end{trapbox}

% ----------------------------------------------------------------------
% PART 7: COMPARATIVE SYNTHESIS & SUMMARY TABLE
% ----------------------------------------------------------------------
\section{Comparative Synthesis \& Summary Table}

Table~\ref{tab:optim_algorithms} compares the primary numerical optimization algorithms used in machine learning.

\begin{table}[htbp]
\centering
\small
\begin{tabularx}{\textwidth}{lXXXXX}
\toprule
\textbf{Algorithm} & \textbf{Order} & \textbf{Per-Iter Cost} & \textbf{Convergence Rate (Convex)} & \textbf{Memory Requirement} & \textbf{Sensitivity to $\kappa$} \\
\midrule
Gradient Descent (GD) & 1st & $\mathcal{O}(d)$ & Sublinear $\mathcal{O}(1/k)$ & $\mathcal{O}(d)$ & Extreme zig-zagging for large $\kappa$ \\
\addlinespace
GD with Momentum & 1st & $\mathcal{O}(d)$ & Accelerated $\mathcal{O}(1/k^2)$ & $\mathcal{O}(d)$ & Damps high-frequency oscillations \\
\addlinespace
Newton-Raphson & 2nd & $\mathcal{O}(d^3)$ & Quadratic $\mathcal{O}(e^{-2^k})$ & $\mathcal{O}(d^2)$ & Completely invariant to condition number \\
\addlinespace
Quasi-Newton (BFGS) & 1st + approx 2nd & $\mathcal{O}(d^2)$ & Superlinear & $\mathcal{O}(d^2)$ & Robust; builds Hessian estimate iteratively \\
\addlinespace
L-BFGS & Limited memory & $\mathcal{O}(m d)$ & Superlinear & $\mathcal{O}(m d)$ ($m \approx 10$) & Industrial standard for medium-scale convex ML \\
\bottomrule
\end{tabularx}
\caption{Comparative Synthesis: Numerical Optimization Algorithms in Machine Learning.}
\label{tab:optim_algorithms}
\end{table}

% ----------------------------------------------------------------------
% PART 8: PYTHON / SCIENTIFIC COMPUTING LAB
% ----------------------------------------------------------------------
\section{Python / Scientific Computing Lab}

The following production script implements Gradient Descent, empirically demonstrates divergence when $\eta > 2/L$, and contrasts isotropic vs.\ ill-conditioned convergence dynamics.

\begin{lstlisting}[language=Python, caption={Gradient Descent Dynamics and Stability Threshold Simulation in Python.}]
import numpy as np

def run_gradient_descent(A: np.ndarray, b: np.ndarray, w0: np.ndarray, lr: float, n_iters: int):
    """Simulates GD on quadratic objective f(w) = 0.5 * w^T A w - b^T w."""
    w = w0.copy()
    history = []
    for k in range(n_iters):
        grad = A @ w - b
        loss = 0.5 * float(w.T @ A @ w) - float(b.T @ w)
        history.append(loss)
        w = w - lr * grad
        if np.isnan(loss) or np.isinf(loss) or loss > 1e12:
            break
    return w, history

def main():
    # Construct an ill-conditioned 2D quadratic: lambda_min = 1.0, lambda_max = 50.0 (kappa = 50)
    L = 50.0
    mu = 1.0
    A = np.array([[L, 0.0], [0.0, mu]])
    b = np.array([0.0, 0.0])  # Minimum at w* = [0, 0]^T
    w0 = np.array([2.0, 2.0])

    stable_lr_bound = 2.0 / L
    optimal_lr = 2.0 / (L + mu)

    print(f"Theoretical Stability Bound (2/L): {stable_lr_bound:.4f}")
    print(f"Optimal Learning Rate 2/(L+mu):   {optimal_lr:.4f}")

    # Case 1: Optimal Learning Rate
    _, hist_opt = run_gradient_descent(A, b, w0, lr=optimal_lr, n_iters=25)
    # Case 2: Sub-optimal Small Learning Rate
    _, hist_small = run_gradient_descent(A, b, w0, lr=0.005, n_iters=25)
    # Case 3: Divergent Learning Rate (eta > 2/L)
    _, hist_div = run_gradient_descent(A, b, w0, lr=0.045, n_iters=10)

    print("\nIteration Loss Comparison:")
    print(f"{'Iter':<6}{'Optimal LR':<18}{'Small LR':<18}{'Divergent LR (> 2/L)':<25}")
    print("-" * 65)
    for i in range(min(10, len(hist_opt))):
        div_val = f"{hist_div[i]:.2e}" if i < len(hist_div) else "OVERFLOW"
        print(f"{i:<6}{hist_opt[i]:<18.4e}{hist_small[i]:<18.4e}{div_val:<25}")

if __name__ == "__main__":
    main()
\end{lstlisting}

% ----------------------------------------------------------------------
% PART 9: COMPREHENSIVE PRACTICE PROBLEMS
% ----------------------------------------------------------------------
\section{Comprehensive Practice Problems}

\subsection*{Part A: Analytical \& Mathematical Problems}

\begin{enumerate}[label=\textbf{P\arabic*.}]
    \item \textbf{Analytical Step-Size Bound on a 2D Quadratic:}
    Let $f(x_1, x_2) = 10 x_1^2 + x_2^2 - 4x_1 - 2x_2$.
    \begin{enumerate}[label=(\alph*)]
        \item Compute the gradient $\nabla f(\mathbf{x})$ and Hessian $\nabla^2 f(\mathbf{x})$.
        \item Find the unique global minimizer $\mathbf{x}^*$.
        \item Find the Lipschitz constant $L$ and strong convexity constant $\mu$.
        \item Calculate the maximum stable step size $\eta_{\max}$ and the optimal step size $\eta^*$.
    \end{enumerate}

    \item \textbf{Derivation of Convergence Rate for Convex $L$-Smooth Functions:}
    Let $f$ be convex and $L$-smooth. Suppose we run Gradient Descent with $\eta = \frac{1}{L}$ for $T$ iterations.
    Using the descent lemma and convexity inequality $f(\mathbf{w}_k) - f^* \le \nabla f(\mathbf{w}_k)^T (\mathbf{w}_k - \mathbf{w}^*)$, prove that:
    \begin{equation}
    f(\mathbf{w}_T) - f^* \le \frac{L \|\mathbf{w}_0 - \mathbf{w}^*\|_2^2}{2T}
    \end{equation}

    \item \textbf{Exact Newton Step on Non-Quadratic Function:}
    Let $f(x) = x^4 + 2x^2 + 4x$.
    \begin{enumerate}[label=(\alph*)]
        \item Compute $f'(x)$ and $f''(x)$.
        \item Starting from initial guess $x_0 = 1$, compute the first two Newton-Raphson iterates $x_1$ and $x_2$.
        \item Contrast the rapid convergence of Newton's method against Gradient Descent with $\eta = 0.05$.
    \end{enumerate}

    \item \textbf{Momentum Acceleration Dynamics:}
    The Heavy-Ball Momentum method updates parameters via:
    \begin{equation}
    \mathbf{w}_{k+1} = \mathbf{w}_k - \alpha \nabla f(\mathbf{w}_k) + \beta (\mathbf{w}_k - \mathbf{w}_{k-1})
    \end{equation}
    Show that for a 1D quadratic $f(w) = \frac{1}{2} a w^2$, this recurrence can be written as a $2 \times 2$ discrete dynamical system. Determine the eigenvalues of the transition matrix and state the conditions on $\alpha, \beta$ for asymptotic stability.

    \item \textbf{Strongly Convex Linear Convergence Rate:}
    Let $f$ be $\mu$-strongly convex and $L$-smooth. Prove that with step size $\eta = \frac{2}{L + \mu}$:
    \begin{equation}
    \|\mathbf{w}_k - \mathbf{w}^*\|_2 \le \left( \frac{\kappa - 1}{\kappa + 1} \right)^k \|\mathbf{w}_0 - \mathbf{w}^*\|_2
    \end{equation}
    where $\kappa = L / \mu$.
\end{enumerate}

\subsection*{Part B: Diagnostic Multiple-Choice Questions}

\begin{enumerate}[label=\textbf{MCQ\arabic*.}]
    \item A loss function $f(\mathbf{w})$ is $L$-smooth with $L = 10$. What is the largest constant step size $\eta$ that guarantees the gradient descent sequence does NOT diverge on an arbitrary convex quadratic loss?
    \begin{enumerate}[label=(\alph*)]
        \item $\eta = 0.1$
        \item $\eta = 0.2$
        \item $\eta = 0.05$
        \item $\eta = 1.0$
    \end{enumerate}

    \item Why does Gradient Descent exhibit slow, zig-zagging convergence on an objective with condition number $\kappa = 10{,}000$?
    \begin{enumerate}[label=(\alph*)]
        \item Because the gradient vector is dominated by the direction of highest curvature, forcing very small step sizes to prevent divergence along that axis, while progress along low-curvature directions is sluggish.
        \item Because the objective function is non-convex.
        \item Because the gradient norm evaluates to zero at all points along the trajectory.
        \item Because the Hessian matrix is singular and cannot be inverted.
    \end{enumerate}
\end{enumerate}
